\DocumentMetadata{
  pdfversion=2.0,
  pdfstandard=ua-2,
  lang=en-US,
  tagging=on,
  tagging-setup={math/setup=mathml-SE,tabsorder=structure}
}
\documentclass[11pt,letterpaper]{article}
\usepackage[margin=0.68in,includefoot]{geometry}
\usepackage{newcomputermodern}
\usepackage{amsmath,amscd,mathtools}
\usepackage{tikz,adjustbox}
\providecommand{\square}{\mdlgwhtsquare}
\usepackage[hidelinks]{hyperref}
\hypersetup{
  pdftitle={Combinatorics PhD Exam, May 2017},
  pdfauthor={University of Florida Department of Mathematics},
  pdfsubject={Graduate mathematics examination},
  pdfdisplaydoctitle=true,
  bookmarksopen=true
}
\setcounter{secnumdepth}{0}
\setlength{\parindent}{0pt}
\setlength{\parskip}{0.28em}
\setlength{\emergencystretch}{3em}
\setlength{\leftmargini}{2.15em}
\setlength{\leftmarginii}{2.65em}
\setlength{\leftmarginiii}{3em}
\pagestyle{empty}
\raggedbottom
\allowdisplaybreaks
\NewTaggingSocketPlug{block/list/label}{examproblem}{%
  \tagstructbegin{tag=Lbl}\tagstructend%
  \tagstructbegin{tag=\LBody}%
  \tagstructbegin{tag=\ProblemHeadingTag,title-o={\ProblemHeadingText}}%
  \tagmcbegin{tag=\ProblemHeadingTag,actualtext={\ProblemHeadingText}}%
  #2%
  \tagmcend%
  \tagstructend%
}
\newenvironment{examproblems}[2]{%
  \def\ProblemHeadingTag{#2}%
  \AssignTaggingSocketPlug{block/list/label}{examproblem}%
  \begin{enumerate}%
  \setcounter{enumi}{#1}%
  \renewcommand{\labelenumi}{\textbf{\theenumi.}}%
  \setlength{\topsep}{0.35em}%
  \setlength{\partopsep}{0pt}%
  \setlength{\itemsep}{0.48em}%
  \setlength{\parsep}{0.18em}%
}{\end{enumerate}}
\newenvironment{examsubparts}{%
  \AssignTaggingSocketPlug{block/list/label}{default}%
  \begin{enumerate}%
  \setlength{\topsep}{0.18em}%
  \setlength{\partopsep}{0pt}%
  \setlength{\itemsep}{0.16em}%
  \setlength{\parsep}{0.1em}%
}{\end{enumerate}}
\newenvironment{examinstructions}{%
  \par\begingroup\itshape
}{%
  \par\endgroup\vspace{0.18em}
}
\makeatletter
\renewcommand\subsection{\@startsection{subsection}{2}{0pt}%
  {-1.25ex plus -.25ex minus -.1ex}%
  {0.55ex plus .1ex}%
  {\normalfont\large\bfseries}}
\renewcommand{\@maketitle}{%
  \newpage\null\vskip -1em%
  {\footnotesize\bfseries
    \href{https://www.ufl.edu/}{University of Florida}, \href{https://math.ufl.edu/}{Department of Mathematics}\par}%
  \vskip -0.5em%
  \tagstructbegin{tag=H1,title={Combinatorics PhD Exam, May 2017}}%
  \tagpdfparaOff%
  \tagmcbegin{tag=H1}%
    {\LARGE\bfseries\href{https://gma.math.ufl.edu/exams/combinatorics-phd/}{Combinatorics PhD Exam}, May 2017\par}%
  \tagmcend%
  \tagpdfparaOn%
  \tagstructend%
  \vskip 0.25em%
  \tagmcbegin{artifact}%
    \hrule height 1pt%
  \tagmcend%
  \vskip 1em%
}
\makeatother
\title{Combinatorics PhD Exam, May 2017}
\author{}
\date{}
\begin{document}
\maketitle
\thispagestyle{empty}
\begin{examproblems}{0}{H2}
\def\ProblemHeadingText{Problem 1.}
\item
Give a combinatorial proof that
\[
k\binom{n}{k}=n\binom{n-1}{k-1}.
\]
\def\ProblemHeadingText{Problem 2.}
\item
You can use results from your course to prove:
\begin{examsubparts}
\item[\textnormal{(a)}]
If a finite partially ordered set has \(mn+1\) elements, then there is either a chain of size \(m+1\) or an antichain of size \(n+1\).
\item[\textnormal{(b)}]
If~\(G\) is a regular bipartite graph of degree at least 2, then~\(G\) has a spanning subgraph that is the union of vertex disjoint cycles.
\end{examsubparts}
\def\ProblemHeadingText{Problem 3.}
\item
Let \(a(n)\) be the number of ways to tile a \(2\times n\) rectangle with \(1\times 1\) squares and \(2\times k\) rectangles for any integer \(k\geq 1\), where a \(2\times 1\) rectangle must be vertical (i.e. no \(1\times 2\) rectangle). For instance, \(a(2)=5\), given by
\par\smallskip
\begin{center}
\begin{adjustbox}{max width=.8\linewidth,max totalheight=6\baselineskip}
\begin{tikzpicture}[
  alt={Five tilings of a 2x2 rectangle, shown left to right: four 1x1 squares; two 1x1 squares on the left with one vertical 2x1 rectangle on the right; one vertical 2x1 rectangle on the left with two 1x1 squares on the right; two vertical 2x1 rectangles; and one 2x2 rectangle.},
  line width=0.8pt,
  line cap=round,
  line join=round
]
% 1. Four 1-by-1 squares
\begin{scope}[xshift=0cm]
    \draw (0,0) rectangle (2,2);
    \draw (1,0) -- (1,2);
    \draw (0,1) -- (2,1);
\end{scope}

% 2. Two 1-by-1 squares on the left and one vertical 2-by-1 tile
\begin{scope}[xshift=3.2cm]
    \draw (0,0) rectangle (2,2);
    \draw (1,0) -- (1,2);
    \draw (0,1) -- (1,1);
\end{scope}

% 3. One vertical 2-by-1 tile and two 1-by-1 squares on the right
\begin{scope}[xshift=6.4cm]
    \draw (0,0) rectangle (2,2);
    \draw (1,0) -- (1,2);
    \draw (1,1) -- (2,1);
\end{scope}

% 4. Two vertical 2-by-1 tiles
\begin{scope}[xshift=9.6cm]
    \draw (0,0) rectangle (2,2);
    \draw (1,0) -- (1,2);
\end{scope}

% 5. One 2-by-2 tile
\begin{scope}[xshift=12.8cm]
    \draw (0,0) rectangle (2,2);
\end{scope}
\end{tikzpicture}
\end{adjustbox}
\end{center}
\smallskip
Set \(a(0)=1\), and find a simple expression for the generating function
\[
A(x)=\sum_{n\geq 0}a(n)x^n.
\]
\def\ProblemHeadingText{Problem 4.}
\item
Let~\(C\) be a linear binary error correcting code. Prove that either every word in~\(C\) has even weight or exactly half the words in~\(C\) have even weight.

\par
Is a similar statement true for ternary codes?
\def\ProblemHeadingText{Problem 5.}
\item
Let~\(\mathcal{L}\) be a regular language. Prove that the language
\[
\mathcal{L}_2=\{w:ww\in\mathcal{L}\}
\]
 is regular.
\def\ProblemHeadingText{Problem 6.}
\item
Prove that the number of \(k\cdots 21\)-avoiding permutations is at most \((k-1)^{2n}\).
\def\ProblemHeadingText{Problem 7.}
\item
Let \(c(n,k)\) be the number of permutations of length~\(n\) with~\(k\) cycles. Let \(a_n\) be the total number of cycles in all permutations of length~\(n\). Prove that
\[
c(n+1,2)=a_n.
\]
\def\ProblemHeadingText{Problem 8.}
\item
Let \(b_n\) be the number of permutations of length~\(n\) whose first ascent is in an even position. (For the decreasing permutation of length~\(n\), we say that its first ascent is in position~\(n\).)

\par
Find an explicit formula for \(b_n\). Your formula can contain one summation sign.
\end{examproblems}
\end{document}
